%5. Global governance: majority support for global democracy on world issues

%===============================================================================================
% - International governance needed to achieve int'l redistribution: are people ready?
%===============================================================================================

International redistribution requires an institution capable of setting common rules, mobilizing resources and implementing collective decisions across countries. The key question is therefore whether citizens are willing to grant international institutions enough authority to act beyond national borders, including for redistributive purposes. 

%===============================================================================================
% - works of Farsan Ghassim
%===============================================================================================
%Fabre 25 NHB, p 1584 
\citet{ghassim_who_2024} find that support for world government strongly depends on its institutional design. Their survey covers 17 countries: Argentina, Australia, Canada, China, Colombia, Egypt, France, Hungary, India, Indonesia, Kenya, Russia, South Korea, Spain, Turkey, the United Kingdom and the United States. Using population weights, support reaches 58\% for an unspecified world government, 71\% for a democratic one, 69\% for one restricted to global issues, and 72\% for an institution combining both characteristics\footnote{Responses were recorded on a six-point scale: ``strongly oppose'', ``oppose'', ``somewhat oppose'', ``somewhat support'', ``support'' and ``strongly support''. There was no neutral midpoint. Here, ``support'' includes the answers ``somewhat support'', ``support'' and ``strongly support''.} \citep[Figure 2]{ghassim_who_2024}.

The authors then focus on the fully specified proposal: a democratic world government empowered to address ``climate change, world poverty, and international peace'', while national governments retain authority over non-global issues. With population weights, support reaches 72\% in the full sample, 75\% in partly free or non-free countries but 56\% in free countries. More generally, support is higher in more populous, less developed, less free and less powerful countries \citep[Figures 2 \& 5]{ghassim_who_2024}.% AF: n'y a-t-il pas la stat pour le monde entier ? encore une fois, on n'utilise que weighted. Aussi, quels pays n'ont pas de majorité pour ça ? %LB: il n'y a pas de stat pour le monde entier. L'annexe en ligne ne préciser pas les stats pays par pays.

The fully specified proposal receives majority support in 16 of the 17 surveyed countries. In these 16 countries, support ranges from 56\% to 82\%, indicating that approval is not confined to a particular region, income group or political regime. The United States constitutes the only exception, with 55\% of respondents opposing the proposal \citep[Figure 6]{ghassim_who_2024}. Overall, these results document broad cross-country support for a democratic global institution whose authority is limited to transnational issues.

Support for global democracy is sufficiently salient to affect stated voting intentions. Across Brazil, Germany, Japan, the United Kingdom and the United States, between 55\% and 74\% of respondents support a directly elected global parliament \citep[Figure 7]{ghassim_global_democracy_2020}. In Brazil and Germany, informing respondents about parties' supposed positions on global democracy shifts aggregate voting intentions from parties described as opposing the proposal toward those described as supporting it by 8 and 12 percentage points, respectively \citep[Figures 144--145]{ghassim_global_democracy_2020}. In Germany, when respondents are told that only the Greens and Die Linke support global democracy, these parties gain 9 and 3 percentage points, respectively, while the SPD and the CDU-CSU each lose 6 percentage points \citep[Figure 145]{ghassim_global_democracy_2020}.

\citet{ghassim_public_2022} ask respondents in Argentina, China, India, Russia, Spain and the United States which institutional features they prefer for a stronger and more representative UN. Making UN decisions binding on every country in areas such as security, environmental and economic matters receives a marginal mean of approximately 0.53\footnote{A marginal mean is the average probability that a profile containing a given institutional feature is selected.}. 
By contrast, limiting their binding force to countries that ``voluntarily accept'' them is the least preferred option, with a marginal mean of about 0.46. The option of having no enforcement mechanism also receives a relatively low marginal mean of 0.48. Collective enforcement by member states and direct enforcement by a strengthened UN reach approximately 0.51 and 0.52, respectively, although the difference between these two options is not statistically significant. Finally, creating a second chamber composed of representatives directly elected by citizens receives a marginal mean of about 0.53, % AF: is there majority support in every country for elections? What about China? %LB: China is not mentioned, and there is no majority in every countries.
while assigning all important decisions to a body in which every member country is represented reaches approximately 0.50 \citep[Figure 1]{ghassim_public_2022}.

%===============================================================================================
% - original results: intl_governance, intl_policy (cf. https://github.com/bixiou/budget_26/tree/main/figures)
%===============================================================================================

\begin{figure}[htbp]
    \centering
    \includegraphics[width=\linewidth]{figures/intl_policy_en.pdf}
    \caption{Evaluations of international policy proposals among respondents in France in 2026.}
    \label{fig:intl_policy}
\end{figure}

\begin{figure}[htbp]
    \centering
    \includegraphics[width=\linewidth]{figures/intl_governance_en.pdf}
    \caption{Preferred mechanisms for making decisions on global issues among respondents in France in 2026.}
    \label{fig:intl_governance}
\end{figure}

% AF: il faut afficher le graphique ici vu qu'il est original
Fabre reports original evidence from an online survey conducted in 2026 among 1,000 respondents representative of the French population. Overall, 60\% favor involving citizens worldwide in decisions on global issues. Among the proposed mechanisms, 75\% evaluate direct worldwide referendums favorably, while 70\% support representative surveys of global public opinion. Citizen assemblies selected by lot, whose decisions would subsequently be implemented by national parliaments, receive 61\% favorable evaluations. By comparison, 54\% favor a directly elected world parliament responsible for most decisions \url{https://github.com/bixiou/budget_26/blob/main/figures/intl_policy_en.pdf} \url{https://github.com/bixiou/budget_26/blob/main/figures/intl_governance_en.pdf}.

